% -----------------------------------------------------------------------------
% Photon observables beyond the differential Compton formula: strict SI
% -----------------------------------------------------------------------------

Klein--Nishina 角分布已经回答了“光子被散射到哪个方向”的问题；实验中还常需要总散射概率，因此下一步是把角分布对全部末态方向积分。这个积分并不引入新的 QED 动力学，它只检验前面振幅、反冲运动学和二体相空间是否能够在两个极限中同时闭合：低能时恢复 Thomson 截面，高能时则出现明显的反冲抑制。

沿用电子静止系的无量纲入射能量
\begin{equation*}
 x\equiv\frac{\hbar\omega}{\me c^2},
 \qquad
 r_e\equiv\frac{\alpha\hbar}{\me c},
\end{equation*}
其中 $r_e$ 是经典电子半径。对\eqnrefs{eq:klein-nishina-SI} 作完整立体角积分，得到
\begin{equation}
 \boxed{
 \begin{aligned}
 \sigma_{\rm KN}(x)=2\pi r_e^2\bigg[&
 \frac{1+x}{x^3}
 \left(
 \frac{2x(1+x)}{1+2x}-\ln(1+2x)
 \right)\\
 &+\frac{\ln(1+2x)}{2x}
 -\frac{1+3x}{(1+2x)^2}
 \bigg].
 \end{aligned}}
 \label{eq:KN-total-r6}
\end{equation}
这个闭式把两个物理尺度直接放在一起。$x\ll1$ 时，电子反冲能量远小于静能，
\begin{equation}
 \sigma_{\rm KN}
 =\frac{8\pi}{3}r_e^2
 \left[1-2x+\order(x^2)\right],
 \label{eq:KN-low-energy-dp49}
\end{equation}
首项正是 Thomson 总截面；$x\gg1$ 时，
\begin{equation}
 \sigma_{\rm KN}
 \simeq
 \frac{\pi r_e^2}{x}
 \left[\ln(2x)+\frac12\right],
 \label{eq:KN-high-energy-dp49}
\end{equation}
说明高能光子虽然仍可发生 Compton 散射，但总截面因相对论反冲按 $\ln x/x$ 下降。低能与高能表达式分别是同一个 Klein--Nishina 结果在 $x\ll1$ 与 $x\gg1$ 两个尺度区间中的渐近展开。

把 Compton 振幅的一条外部光子和电子线作交叉延拓，会得到更强的检验：两个真实光子能否把能量转化成一对真实电子和正电子。考虑
\begin{equation*}
 \gamma(k_1,\lambda_1)+\gamma(k_2,\lambda_2)
 \longrightarrow e^-(p_1,s_1)+e^+(p_2,s_2),
\end{equation*}
并继续把 $k_i^\mu$、$p_i^\mu$ 作为物理四动量。两个入射光子满足 $k_1^2=k_2^2=0$，末态满足 $p_1^2=p_2^2=\me^2c^2$。定义
\begin{equation}
 s\equiv(k_1+k_2)^2,
 \qquad
 \beta_{\rm pair}\equiv
 \sqrt{1-\frac{4\me^2c^2}{s}},
 \label{eq:BW-beta-pair-dp49}
\end{equation}
另外定义两种内部电子连接对应的四动量不变量
\begin{equation}
 t\equiv(k_1-p_1)^2,
 \qquad
 u\equiv(k_1-p_2)^2,
 \qquad
 s+t+u=2\me^2c^2.
 \label{eq:BW-stu-si-r12}
\end{equation}
最后一个关系只使用四动量守恒和四条外腿的质量壳条件；它使两张树图的两个电子传播子分母可以统一写成 $t-\me^2c^2$ 与 $u-\me^2c^2$。其中 $c\beta_{\rm pair}$ 是质心系末态电子和正电子的速率。真实二体末态要求
\begin{equation*}
 s\ge 4\me^2c^2,
\end{equation*}
所以 $\beta_{\rm pair}=0$ 正好是对产生阈值。

最低阶仍只有两个 QED 顶角，但沿内部电子线有两种连接次序。两图必须在振幅层相加。按\eqnrefs{eq:si-reduced-propagators-r10,eq:si-reduced-vertices-r10} 的约化 SI 图规则，树级振幅为
\begin{equation}
 \boxed{
 \mathcal M_{\rm BW}
 =-g_{\rm em}^2\bar u(p_1)\left[
 \slashed\varepsilon_1
 \frac{\slashed p_1-\slashed k_1+\me c}{t-\me^2c^2}
 \slashed\varepsilon_2
 +\slashed\varepsilon_2
 \frac{\slashed p_1-\slashed k_2+\me c}{u-\me^2c^2}
 \slashed\varepsilon_1
 \right]v(p_2).}
 \label{eq:BW-amplitude-r68}
\end{equation}
它与 Compton 振幅属于同一解析四点函数的不同交叉道。自旋--偏振求和和二体相空间随后给出质心系角分布
\begin{equation}
 \boxed{
 \frac{\dd\sigma_{\rm BW}}{\dd\Omega}
 =\frac{r_e^2}{4}\,\beta_{\rm pair}(1-\beta_{\rm pair}^2)
 \frac{
 1-\beta_{\rm pair}^4\cos^4\theta
 +2\beta_{\rm pair}^2(1-\beta_{\rm pair}^2)\sin^2\theta
 }{(1-\beta_{\rm pair}^2\cos^2\theta)^2}.}
 \label{eq:BW-dsigma-si-dp49}
\end{equation}
这里 $\theta$ 是入射光子 $\mathbf k_1$ 与末态电子 $\mathbf p_1$ 的夹角。再对立体角积分，得到 Breit--Wheeler 总截面
\begin{equation}
 \boxed{
 \sigma_{\rm BW}
 =\frac{\pi r_e^2}{2}(1-\beta_{\rm pair}^2)
 \left[
 (3-\beta_{\rm pair}^4)
 \ln\frac{1+\beta_{\rm pair}}{1-\beta_{\rm pair}}
 -2\beta_{\rm pair}(2-\beta_{\rm pair}^2)
 \right].}
 \label{eq:BW-total-si-dp49}
\end{equation}
阈值附近
\begin{equation}
 \sigma_{\rm BW}
 =\pi r_e^2\beta_{\rm pair}
 +\order(\beta_{\rm pair}^3),
 \label{eq:BW-threshold-dp49}
\end{equation}
线性开启来自二体末态相空间；在超相对论极限 $\beta_{\rm pair}\to1$，前因子 $1-\beta_{\rm pair}^2$ 又压低总截面。因此总截面必在阈值与超相对论区之间出现有限峰值。这个模型把“传播子中的虚电子”与“末态中的真实电子”严格区分开：前者只受顶角四动量守恒约束，后者还必须满足 LSZ 选出的真实质量壳。

\begin{figure}[htbp]
 \centering
 \begin{tikzpicture}[scale=.76, baseline=(current bounding box.center)]
\begin{scope}[xshift=-4cm]
\coordinate (v1) at (-.7,0);\coordinate (v2) at (.7,0);
\draw[photon] (-2.2,1.2)--(v1);\draw[photon] (-2.2,-1.2)--(v2);
\draw[fermion] (v1)--(v2);
\draw[fermion] (v1)--(2.2,1.2);\draw[fermion] (2.2,-1.2)--(v2);
\node[above=3pt,fill=white,inner sep=1pt] at (0,0){$p_1-k_1$};
\end{scope}
\node at (0,0){$+$};
\begin{scope}[xshift=4cm]
\coordinate (v1) at (-.7,0);\coordinate (v2) at (.7,0);
\draw[photon] (-2.2,1.2)--(v2);\draw[photon] (-2.2,-1.2)--(v1);
\draw[fermion] (v1)--(v2);
\draw[fermion] (v2)--(2.2,1.2);\draw[fermion] (2.2,-1.2)--(v1);
\node[above=3pt,fill=white,inner sep=1pt] at (0,0){$p_1-k_2$};
\end{scope}
\end{tikzpicture}

 \caption{Breit--Wheeler 对产生的两张电子交换树图。两张图对应内部电子线上的两种顶角次序，必须先相干相加；四条外腿均由渐近单粒子极点固定在真实质量壳。}
 \label{fig:breit-wheeler-two-r12}
\end{figure}

\begin{figure}[htbp]
 \centering
 \includegraphics[width=.72\linewidth]{breit_wheeler_total.png}
 \caption{Breit--Wheeler 总截面的无量纲形状，纵轴为 $\sigma_{\rm BW}/(\pi r_e^2)$。阈值处二体相空间关闭，高能区又受反冲与传播子结构抑制，因此中间能区出现有限最大值。}
 \label{fig:breit-wheeler-total-r6}
\end{figure}

\begin{note}[外腿质量壳与内部传播变量]
真实外部光子由 LSZ 的单粒子极点选定，因此满足 $k^2=0$ 并只保留两个横向偏振。内部光子或电子的四动量则由相邻顶角的四动量守恒确定，是时序 Green 函数的 Fourier 积分变量，一般不满足对应自由粒子的质量壳。对有质量媒介场同样如此：外部真实量子满足自身质量壳，静态交换所取的内部四动量却通常是类空的。真实外态与内部传播变量的差别来自是否施加渐近单粒子极点条件，而不是来自能量守恒是否成立。
\end{note}
